\section{Numerical charge slope of the certified backgrounds}\label{supp:charge-slope}
This section reports a numerical evaluation added in preprint 0.53. It is
not an interval enclosure and does not change any certificate.

\paragraph{Definitions and method.}
With the convention of the main paper, $Q=2\omega\int f^2\dd^3x=
8\pi\omega\int_0^\infty r^2f^2\dd r$ and $G=\int|\nabla f|^2\dd^3x$.
Profiles of Eq.~\eqref{eq:profile} were computed by collocation with a
relative tolerance $10^{-8}$, a Robin condition
$f'+(\kappa_0+1/R)f=0$ at the outer radius
$R=F\max(40,16/\kappa_0,R_{\mathrm{tw}}+30)$ and an analytic exponential
tail correction to $Q$. The slope $D=dQ/d(\omega^2)$ was taken as a
central difference with steps $\delta=2\times10^{-3}$ and $10^{-3}$
($F=1$) and $\delta=10^{-3}$ ($F=1.25$); the noise measure $\eta$ is
the sum of the absolute differences between consecutive variants, and
$dQ/d\omega=2\omega D$. Before the calculation, a background was to be
classified as lying on the branch $dQ/d\omega<0$ only if all three
variants were negative and $|D|>10\eta$.

\begin{table}[htbp]\centering\small
\caption{Charge and charge slope at the midpoints of the certified
parameter boxes (Table~\ref{tab:proof-cases}). $D$ is the variant with
$\delta=10^{-3}$, $F=1$. Numerical values, not enclosures.}
\label{tab:charge-slope}
\begin{tabular}{lllllr}\toprule
Case&$\omega^2$&$Q$&$D=dQ/d(\omega^2)$&$dQ/d\omega$&$|D|/\eta$\\\midrule
C0&$0.797677$&$189.14436$&$-1326.19$&$-2368.9$&$10\,560$\\
C1&$0.685129$&$573.40356$&$-7743.01$&$-12\,818$&$3\,919$\\
C2&$0.754496$&$266.78523$&$-2394.87$&$-4160.5$&$7\,609$\\\bottomrule
\end{tabular}
\end{table}

\paragraph{Charge minimum.}
On a grid $\omega^2=0.59,\dots,0.998$ the charge has exactly one local
minimum. A golden-section refinement gives
$\omega_Q^2\simeq0.927176$ and $Q_{\min}\simeq111.860$, independently of
$F\in\{1,1.25\}$. $Q$ decreases from $\omega^2=0.59$ to $\omega_Q^2$ and
increases above it. The three certified backgrounds lie $0.13$--$0.24$
below $\omega_Q^2$ in $\omega^2$. Two profiles of equal charge reported
in earlier project calculations, at $\omega^2\simeq0.766$ and $0.993$, lie
on opposite sides of this minimum. Below $\omega^2=0.59$ the collocation
routine found no accepted profile; this does not affect the three
backgrounds.

\paragraph{Controls.}
(i)~The two equal-charge profiles were reproduced to relative accuracy
$7\times10^{-8}$ (threshold $10^{-5}$).
(ii)~Cubic nonlinear Schr\"odinger limit. For $\omega\to1$,
$Q\kappa_0/\omega$ must approach $N_R=\int g^2\dd^3s$, where $g$ is the
positive ground state of $g''+2g'/s=g-g^3$. A separate shooting
calculation gives $g(0)\simeq4.33738768$ and $N_R\simeq18.897251$, stable
to $4\times10^{-11}$ and $4\times10^{-12}$ relative under step and
outer-radius changes. The computed $Q\kappa_0/\omega$ is $20.186$,
$19.529$ and $19.147$ at $\omega^2=0.99$, $0.995$ and $0.998$; the
relative excess over $N_R$ falls monotonically to $1.3\,\%$ (threshold
$3\,\%$ at $\omega^2=0.998$).
(iii)~Virial identity. Derrick scaling and the Hellmann--Feynman
relation for the action $\int(|\nabla f|^2+\mathcal U(f^2)-\omega^2f^2)
\dd^3x$ give the exact identity $dG/d(\omega^2)=-3Q/(4\omega)$. With
central differences of $G$ at $\delta=5\times10^{-4}$ and
$2.5\times10^{-4}$, the largest relative deviation is approximately
$1.2\times10^{-5}$ and $3.0\times10^{-6}$ respectively (threshold
$10^{-4}$); halving $\delta$ reduces it by a factor of four, as expected
for the truncation error of the difference quotient.
(iv)~An independent shooting calculation with a fourth-order
Runge--Kutta integrator, an exact Taylor start near the origin and an
analytic Yukawa tail reproduced $Q$ at eleven comparison points to
$3\times10^{-10}$ relative, the slopes $D$ at the same finite-difference
step $\delta=10^{-3}$ to six or more digits and $\omega_Q^2$ to
$10^{-6}$; with steps down to $\delta=2.5\times10^{-4}$ it gives
$|D|/\eta>10^4$ at all three backgrounds. Between $\delta=10^{-3}$ and
$2.5\times10^{-4}$ the slopes change by at most $8\times10^{-5}$
relative (truncation error of the difference quotient); the signs are
unaffected. The displayed digits of $\omega_Q^2$ and $Q_{\min}$ are
numerical values, not a validated accuracy. It also obtained
accepted profiles down to $\omega^2=0.511$, where $Q$ approaches the
leading thin-wall charge from above. In that calculation an additional
classification check of the shooting bracket, not part of the
prespecified acceptance rule, failed at the level of floating-point
rounding; the values above are unaffected.

\paragraph{Scope.}
For the $l=0$ ansatz with a real growth rate, Kovtun, Nugaev and
Shkerin arrive at the Vakhitov--Kolokolov criterion $dQ/d\omega>0$ for
instability, in the same sign convention~\cite{Kovtun2018,VakhitovKolokolov1973}.
Provided the spectral hypotheses of this radial criterion hold for the
nodeless sextic branch considered here, the negative slopes give
numerical evidence against a purely growing radial instability of the
three backgrounds; these hypotheses are not independently verified by
the charge-slope calculation. The slopes do not address
oscillatory instabilities, the sectors $l\ge1$, orbital or nonlinear
stability, the persistence of the embedded modes at finite amplitude or
a lifetime. Programs and result records of both calculations are
included in archive revision~3 (directory \texttt{charge-slope/}).
